% 原书物理页=印刷页28--30；1.8 全节.
\section{Lie 群在流形上的作用}\label{sec:lie-1-8}

这里想要表达的想法是：光滑流形 $M$ 的微分同胚群表现得就像一个 Lie 群，其 Lie 代数是 $M$ 上向量场全体构成的代数 $\mathcal V(M)$，配备由导子引出的 Lie 括号. 于是，Lie 群 $G$ 在 $M$ 上的一个光滑作用可以被``求导''成 $\operatorname{Lie}(G)$ 在 $M$ 上函数的一个作用，而我们也可以把它表述为一个 Lie 代数态射 $\operatorname{Lie}(G)\longrightarrow\mathcal V(M)$.

\begin{numpara}[向量场的 Lie 代数]\label{np:lie-1-8-1}
设 $M$ 是一个光滑流形. 一个向量场 $\xi\in\mathcal V(M)$ 确定光滑函数的 $\mathbb R$-代数 $\mathcal C^{\infty}(M)$ 的一个导子 $\partial_{\xi}$，从而给出一个线性映射
\[
\begin{array}{rcl}
\mathcal V(M) & \longrightarrow & \operatorname{Der}_{\mathbb R}\bigl(\mathcal C^{\infty}(M)\bigr)\\[2pt]
\xi & \longmapsto & \partial_{\xi}:\ \bigl(m\mapsto d_{m}f(\xi(m))\bigr)\ .
\end{array}
\]
我们知道这个映射是\emph{双射}的（参见微分几何课程）. 由此在 $\mathcal V(M)$ 上得到一个 Lie 括号 $(\xi,\xi')\mapsto[\xi,\xi']$，它由等式 $\partial_{[\xi,\xi']}=[\partial_{\xi},\partial_{\xi'}]=\partial_{\xi}\partial_{\xi'}-\partial_{\xi'}\partial_{\xi}$ 刻画. 我们还有表达式：
\[
(*)\quad [\xi,\xi']=\frac{d}{dt}\Bigl((\varphi^{\xi}_{t})^{*}\xi'\Bigr)_{|t=0}:\ m\mapsto\frac{d}{dt}\Bigl((d_{m}\varphi^{\xi}_{t})^{-1}\Bigl(\xi'\bigl(\varphi^{\xi}_{t}(m)\bigr)\Bigr)\Bigr)_{|t=0}
\]
\end{numpara}

\begin{example*}
在 $M=\mathbb R^{n}$ 中，有 $[x^{i}\frac{\partial}{\partial x_{j}},\,x^{k}\frac{\partial}{\partial x_{l}}]=\delta_{jk}x^{i}\frac{\partial}{\partial x_{l}}-\delta_{il}x^{k}\frac{\partial}{\partial x_{j}}$.
\end{example*}

\begin{numpara}[光滑作用的导数]\label{np:lie-1-8-2}
现在假设一个 Lie 群 $G$ 光滑地作用在 $M$ 上，也就是说，给出该作用的态射 $G\times M\longrightarrow M,\ g\mapsto gm$ 是光滑的. 把与之相伴的群态射记作 $\alpha:G\longrightarrow\operatorname{Diff}(M)$. 于是可以把这个作用``求导''，得到
\[
\begin{array}{rcl}
\text{``}d\alpha\text{''}:\quad\operatorname{Lie}(G) & \longrightarrow & \mathcal V(M)\\[2pt]
X & \longmapsto & \text{``}d\alpha\text{''}(X):\ m\mapsto d_{e}(g\mapsto gm)(X)
\end{array}
\]
与向量场 $d\alpha(X)$ 相伴的导子 $\partial_{d\alpha(X)}$ 则由下式给出：
\[
\partial_{d\alpha(X)}(f):\ \bigl(m\mapsto d_{m}f\bigl(d_{e}(g\mapsto gm)(X)\bigr)=d_{e}(g\mapsto f(gm))(X)\bigr),\ \forall f\in\mathcal C^{\infty}(M)
\]
\end{numpara}

\begin{example*}
考察 $G=\mathrm{GL}_{n}(\mathbb R)$ 在 $M=\mathbb R^{n}$ 上的作用. $\mathfrak{gl}_{n}(\mathbb R)$ 有一组由初等矩阵构成的自然基 $(E_{ij})_{1\leqslant i,j\leqslant n}$. 作为 $G$ 在 $e$ 处的切向量，它们给出一些点态导子，后者可以简单地写为
\[
d_{e}f(E_{ij})=\left(\frac{\partial}{\partial x_{ij}}\right)_{|(\delta_{ij})_{i,j}}\left((x_{ij})_{i,j}\mapsto f\Bigl(\sum_{ij}x_{ij}E_{ij}\Bigr)\right),\ \forall f\in\mathcal C^{\infty}(G).
\]
\end{example*}

由此，对 $x=(x_{i})_{i}\in\mathbb R^{n}$ 有
\[
\partial_{d\alpha(E_{ij})}(f)(x)=\left(\frac{\partial}{\partial g_{ij}}\bigl((g_{ij})_{ij}\mapsto f\bigl((g_{ij}).x\bigr)\bigr)\right)_{|(\delta_{ij})_{i,j}}=x_{j}\frac{\partial f}{\partial x_{i}}(x),
\]
从而 $\partial_{d\alpha(E_{ij})}=x_{j}\frac{\partial}{\partial x_{i}}$. 可以注意到，$E_{ij}\mapsto x_{j}\frac{\partial}{\partial x_{i}}$ 仅在\emph{相差一个符号}的意义下与括号相容. 下面的命题将澄清这一状况.

\begin{example*}
考察 $\mathrm{SO}(3,\mathbb R)$ 在 $\mathbb R^{3}$ 上的作用. 把 $\mathbb R^{3}$ 中的坐标记作 $x,y,z$. 容易验证 $\mathfrak{so}(3,\mathbb R)$ 有一组由三个矩阵 $X:=E_{32}-E_{23}$、$Y:=E_{13}-E_{31}$ 与 $Z=E_{21}-E_{12}$ 构成的基. 于是前面的公式给出
\[
\partial_{d\alpha(X)}=y\frac{\partial}{\partial z}-z\frac{\partial}{\partial y},\quad \partial_{d\alpha(Y)}=z\frac{\partial}{\partial x}-x\frac{\partial}{\partial z},\quad \partial_{d\alpha(Z)}=x\frac{\partial}{\partial y}-y\frac{\partial}{\partial x}.
\]
我们注意到，所有这些向量场都与球面 $S^{2}$ 相切.
\end{example*}

\begin{proposition*}
采用上面的记号，``$d\alpha$''是一个 $\mathbb R$-Lie 代数的\newterm{lie-anti-morphism}{反态射}{anti-morphism}，也就是说 $[d\alpha(X),d\alpha(Y)]=-d\alpha([X,Y])$ 对所有 $X,Y\in\operatorname{Lie}(G)$ 成立.
\end{proposition*}

\begin{proof}
i) 可以验证 $d\alpha(X)$ 的流由 $(t,m)\mapsto\exp(tX)m$ 给出. 只需求导并加以确认.
ii) 对一切 $g\in G$，有 $d\alpha(\operatorname{Ad}_{g}(X))=g_{*}(d\alpha(X))$. 事实上，对一切 $m\in M$：
\begin{align*}
d\alpha(\operatorname{Ad}_{g}(X))(m) &= d_{e}(h\mapsto hm)(\operatorname{Ad}_{g}(X))=d_{e}(h\mapsto ghg^{-1}m)(X)\\
&= d_{g^{-1}m}\alpha_{g}\circ d_{e}(h\mapsto hg^{-1}m)(X)=d_{g^{-1}m}\alpha_{g}\bigl(d\alpha(X)(g^{-1}m)\bigr)\\
&= \bigl(g_{*}(d\alpha(X))\bigr)(m)
\end{align*}
iii) 在 $t=0$ 处对恒等式
\[
d\alpha(\operatorname{Ad}_{\exp(tX)}(Y))=\exp(tX)_{*}(d\alpha(Y))=\exp(-tX)^{*}(d\alpha(Y)).
\]
求导. 我们得到恒等式 $d\alpha(\operatorname{ad}_{X}(Y))=\frac{d}{dt}\bigr|_{t=0}\bigl(t\mapsto(\varphi^{d\alpha(-X)}_{t})^{*}(d\alpha(Y))\bigr)$，再由向量场括号的表达式 $(*)$，便得 $d\alpha([X,Y])=-[d\alpha(X),d\alpha(Y)]$.
\end{proof}

\begin{numpara}[稳定子群的 Lie 代数]\label{np:lie-1-8-3}
设 $m$ 是 $M$ 的一个元素. 把 $m$ 在 $G$ 中的稳定子群记作 $G_{m}$. 它是 $G$ 的一个闭子群，因此我们知道它是一个 Lie 子群. 前面的命题向我们提示了一种计算其 Lie 代数的方法.
\end{numpara}

\begin{corollary*}
$\operatorname{Lie}(G_{m})=\{X\in\operatorname{Lie}(G),\ \frac{d}{dt}\bigr|_{t=0}(t\mapsto\exp(tX)m)=0\}=\operatorname{Ker}(d\alpha(m))$.
\end{corollary*}

\begin{proof}
我们知道 $\operatorname{Lie}(G_{m})=\{X\in\operatorname{Lie}(G),\ \forall t,\ \exp(tX)\in G_{m}\}=\{X\in\operatorname{Lie}(G),\ \forall t,\ \exp(tX)m=m\}$. 于是 $X\in\operatorname{Lie}(G_{m})\Rightarrow\frac{d}{dt}(t\mapsto\exp(tX)m)=0$.
反之，假设 $\frac{d}{dt}\bigr|_{t=0}(t\mapsto\exp(tX)m)=0$. 这也可以写成 $d_{e}(g\mapsto gm)(X)=d\alpha(X)(m)=0$（即 $m$ 是向量场 $d\alpha(X)$ 的一个零点），这蕴含常值曲线 $t\mapsto m$ 是场 $d\alpha(X)$ 的一条积分曲线.
而我们早已看到，$t\mapsto\exp(tX)m$ 也是它的一条积分曲线（参见上述命题证明中的 i)）. 由唯一性，对一切 $t$ 有 $\exp(tX)m=m$.
第二个等式是显然的.
\end{proof}

\begin{example*}
我们可以重新算出各个经典群的 Lie 代数. 例如，$\mathrm{O}(\Phi)\subset\mathrm{GL}(V)$ 是非退化对称双线性型 $\Phi$ 在 $\mathrm{GL}(V)$ 作用于由所有这些型构成的集合（一个向量空间的开集）时的稳定子群. 我们得到
\begin{align*}
\operatorname{Lie}(\mathrm{O}(\Phi)) &= \{X\in\operatorname{Lie}(G),\ \frac{d}{dt}\bigr|_{t=0}(t\mapsto\Phi(\exp(tX)v,\exp(tX)w)=0,\ \forall v,w\in V\}\\
&= \{X\in\operatorname{Lie}(G),\ \Phi(Xv,w)+\Phi(v,Xw)=0,\ \forall v,w\in V\}\\
\end{align*}
\footnote{译注：原书此式第一行 $(t\mapsto\Phi(\exp(tX)v,\exp(tX)w)$ 的外括号未闭合，照录原文.}
\end{example*}

\begin{corollary*}
采用上面的记号，映射 $g\mapsto gm$ 诱导出一个单的浸入 $G/G_{m}\hookrightarrow M$，其像是 $m$ 在 $G$ 作用下的\newterm{lie-orbit}{轨道}{orbit} $\mathcal O_{m}$. 若作用是\newterm{lie-transitive}{传递的}{transitive}，则由此得到一个微分同胚 $G/G_{m}\xrightarrow{\ \sim\ }M$.
\end{corollary*}

\begin{proof}
由商的构造，映射 $g\mapsto gm$ 确实诱导出一个光滑映射 $G/G_{m}\longrightarrow M$. 由于 $\operatorname{Lie}(G_{m})=\operatorname{Ker}(d\alpha(m))$，诱导的映射 $T_{e}(G/G_{m})=\operatorname{Lie}(G)/\operatorname{Lie}(G_{m})\longrightarrow T_{m}M$ 是单射. 经平移可知它在每一点都是单射，故该映射确是一个浸入.
\end{proof}

\begin{example*}
$\mathrm{SO}(n,\mathbb R)$ 从点 $m=(0,\cdots,0,1)$ 出发在球面 $S^{n}$ 上的作用诱导出一个微分同胚 $\mathrm{SO}(n,\mathbb R)/\mathrm{SO}(n-1,\mathbb R)\xrightarrow{\ \sim\ }S^{n}$，从而给出 $\mathrm{SO}(n,\mathbb R)$ 在 $S^{n}$ 上以 $\mathrm{SO}(n-1,\mathbb R)$ 为纤维的一个纤维化.
\end{example*}

\begin{numpara}[$G$ 上的 Lie 代数与不变向量场]\label{np:lie-1-8-4}
现在假设 $M=G$，并带有右平移作用：$R_{g}(h):=hg^{-1}$ 对一切 $g,h\in G$. 那么
\[
\text{``}dR\text{''}(X)=\Bigl(h\mapsto d_{e}\bigl(g\mapsto hg^{-1}\bigr)(X)=d_{e}L_{h}\circ d_{e}i(X)=-d_{e}L_{h}(X)\Bigr)=-\xi_{X}.
\]
（与前面一样，$\xi_{X}$ 正是与 $X$ 相伴的左不变向量场.）
\end{numpara}

\begin{corollary*}
引理~\ref{np:lie-1-7-5} 给出的同构 $\operatorname{Lie}(G)\xrightarrow{\ \sim\ }\mathcal V(G)^{L}$ 是一个 $\mathbb R$-Lie 代数的同构.
\end{corollary*}

\begin{proof}
事实上，$\xi_{[X,Y]}=-\text{``}dR\text{''}([X,Y])=[\text{``}dR\text{''}(X),\text{``}dR\text{''}(Y)]=[-\xi_{X},-\xi_{Y}]=[\xi_{X},\xi_{Y}]$.
\end{proof}

把 Lie 代数显式地转换成不变向量场或与之相伴的导子常常是有用的，尤其是在物理中. 导子的公式为：
\[
\partial_{X}(f):\ \bigl(x\mapsto d_{x}f(\xi_{X}(x))=d_{x}f\circ d_{e}L_{x}(X)=d_{e}(g\mapsto f(xg))(X)\bigr).
\]
\begin{example*}
考察 $G=\mathrm{GL}_{n}(\mathbb R)$ 的情形，并计算与初等矩阵 $E_{ij}$ 相伴的不变导子 $\partial_{E_{ij}}$. 由此推得，对 $x=(x_{ij})_{i,j}$ 有
\begin{align*}
\partial_{E_{ij}}(f)(x)&=\left(\frac{\partial}{\partial g_{ij}}\right)_{|(\delta_{ij})_{ij}}\left((g_{ij})_{i,j}\mapsto f\Bigl(\bigl(\sum_{ij}x_{ij}E_{ij}\bigr)\bigl(\sum_{ij}g_{ij}E_{ij}\bigr)\Bigr)\right)\\
&=\sum_{k=1}^{n}x_{ki}\frac{\partial f}{\partial x_{kj}}(x)
\end{align*}
从而 $\partial_{E_{ij}}=\sum_{k=1}^{n}x_{kj}\frac{\partial}{\partial x_{ki}}$\footnote{译注：原文如此. 按上一行的计算以及公式 $\partial_{X}(f):x\mapsto d_{x}f(\xi_{X}(x))$，此处右端应为 $\sum_{k=1}^{n}x_{ki}\frac{\partial}{\partial x_{kj}}$，即结论一行中的下标 $i,j$ 疑互换了.}.
\end{example*}
